\documentclass[12pt,a4paper]{article}
\usepackage[utf8]{inputenc}
\usepackage[T1]{fontenc}
\usepackage[spanish]{babel}
\usepackage{geometry}
\geometry{left=2.5cm, right=2.5cm, top=2.5cm, bottom=2.5cm}
\usepackage{amsmath}
\usepackage{booktabs} % Para tablas con acabado profesional
\usepackage{parskip}

\title{\textbf{Análisis de Rendimiento y Prestaciones en Arquitectura del Computador}}
\author{\textbf{Renny Andrade}}
\date{}

\begin{document}
	
	\maketitle
	
	\section{Fundamentos de las Prestaciones}
	
	En el ámbito de la informática, el rendimiento o la \textit{prestación} de un sistema está intrínsecamente ligado a la velocidad con la que procesa la información. Para optimizar las prestaciones de un equipo, el objetivo primordial consiste en reducir al máximo el tiempo requerido para completar una tarea determinada. Matemáticamente, la prestación de una computadora $X$ se define de forma inversa a su tiempo de procesamiento:
	
	\begin{equation}
		\text{Prestaciones}_X = \frac{1}{\text{Tiempo de ejecución}_X}
	\end{equation}
	
	A partir de esta relación, si comparamos dos sistemas distintos ($X$ e $Y$) y determinamos que la computadora $X$ ofrece un rendimiento superior a la computadora $Y$, podemos deducir formalmente la siguiente cadena de desigualdades:
	
	\begin{align}
		\text{Prestaciones}_X &> \text{Prestaciones}_Y \\[0.3cm]
		\frac{1}{\text{Tiempo de ejecución}_X} &> \frac{1}{\text{Tiempo de ejecución}_Y} \\[0.3cm]
		\text{Tiempo de ejecución}_Y &> \text{Tiempo de ejecución}_X
	\end{align}
	
	Este desglose demuestra de manera lógica que un equipo con mayores prestaciones completará el mismo trabajo en un intervalo temporal más corto, convirtiéndose de este modo en una alternativa mucho más eficiente.

    \clearpage
	
	\section{Cuantificación Relativa del Rendimiento}
	
	Para evaluar el impacto de las mejoras de hardware, los ingenieros calculan una relación cuantitativa que permite expresar cuántas veces un sistema es más veloz que otro. Tomando como base el escenario previo, si el computador $X$ aventaja al computador $Y$, es habitual enunciar que ``la velocidad de la máquina $X$ es $n$ veces mayor que la de la máquina $Y$''. Dicho factor multiplicador ($n$) se obtiene mediante el cociente de sus rendimientos:
	
	\begin{equation}
		\frac{\text{Prestaciones}_X}{\text{Prestaciones}_Y} = n
	\end{equation}
	
	Dado que medir prestaciones puras de forma directa es complejo, resulta más práctico derivar este factor utilizando la métrica del tiempo de ejecución. Puesto que el rendimiento es inversamente proporcional al tiempo, la proporción se invierte, quedando establecida la ecuación general de comparación de la siguiente manera:
	
	\begin{equation}
		n = \frac{\text{Prestaciones}_X}{\text{Prestaciones}_Y} = \frac{\text{Tiempo de ejecución}_Y}{\text{Tiempo de ejecución}_X}
	\end{equation}
	
	\section{Factores Determinantes en el Rendimiento de la CPU}
	
	Cuando nos enfocamos exclusivamente en la unidad central de procesamiento (CPU), la medición del rendimiento se centra en el tiempo que le toma al procesador procesar un programa de software. Los diseñadores de hardware emplean métricas basadas en los ciclos del reloj del procesador para evaluar de forma granular cómo impacta una modificación física en la experiencia del usuario. 
	
	La relación elemental que gobierna el tiempo de procesamiento en función de la cantidad de pulsos de reloj y la duración de cada uno es:
	
	\begin{equation}
	\text{Tiempo de ejecución de la CPU} = (\text{Ciclos de reloj del programa}) \times (\text{Tiempo del ciclo de reloj})
\end{equation}
	
	Alternativamente, dado que la frecuencia de reloj (medida en hercios) representa el inverso multiplicativo de la duración de un ciclo individual, la ecuación puede reescribirse en términos de velocidad de frecuencia:
	
	\begin{equation}
		\text{Tiempo de ejecución de la CPU} = \frac{\text{Ciclos de reloj del programa}}{\text{Frecuencia del reloj}}
	\end{equation}

    \clearpage
	
	\subsection{La Ecuación Clásica de Rendimiento}
	
	Para profundizar en el análisis, es indispensable vincular las métricas anteriores con el software que se está ejecutando. Esto se logra desglosando los ciclos de reloj totales en función del número de instrucciones ejecutadas y el promedio de ciclos necesarios para procesar cada una de ellas (métrica conocida como CPI o \textit{Ciclos Por Instrucción}).
	
	La fórmula estructural clásica se define como:
	
	\begin{equation}
		\text{Tiempo de ejecución} = \text{Número de instrucciones} \times \text{CPI} \times \text{Tiempo de ciclo}
	\end{equation}
	
	Bajo el mismo principio físico de la frecuencia, la ecuación de rendimiento se puede consolidar dividiendo el producto del software entre la tasa de ciclos por segundo de la arquitectura:
	
	\begin{equation}
		\text{Tiempo de ejecución} = \frac{\text{Número de instrucciones} \times \text{CPI}}{\text{Frecuencia de reloj}}
	\end{equation}
	
	\subsection{Resumen de Componentes y Unidades métricas}
	
	Para interpretar y aplicar correctamente las ecuaciones de rendimiento, se debe tener claridad sobre los componentes involucrados y las unidades en las que se cuantifican:
	
	\begin{table}[htbp]
		\centering
		\caption{Métricas esenciales en la evaluación de prestaciones de la CPU.}
		\vspace{0.3cm}
		\begin{tabular}{ll}
			\toprule
			\textbf{Componente de las prestaciones} & \textbf{Unidad de medida estándar} \\
			\midrule
			Tiempo de ejecución de CPU & Segundos por programa \\
			Número de instrucciones & Instrucciones ejecutadas por programa \\
			Ciclos por instrucción (CPI) & Número medio de ciclos por instrucción \\
			Tiempo de ciclo del reloj & Segundos por ciclo de reloj \\
			\bottomrule
		\end{tabular}
	\end{table}
	
\end{document}